\chapter{Metodología y entorno de desarrollo}

Este capítulo detalla el marco metodológico, las herramientas técnicas y los procedimientos de seguridad empleados para la consecución de los objetivos. Se describe un enfoque centrado en la reproducibilidad y la eficiencia bajo restricciones de recursos.

\section{Gestión del ciclo de vida del proyecto}

El desarrollo de este Trabajo de Fin de Máster se ha articulado bajo una metodología iterativa e incremental, permitiendo ajustar el sistema a medida que se descubrían las particularidades de las gramáticas mayas y las limitaciones de la API de Gemini.

\subsection{Planificación bajo restricciones}

Siguiendo las directrices de la tutoría académica, el proyecto se diseñó para ser ejecutado en una carga de trabajo acotada de 250 horas. Esta restricción temporal ha sido un factor determinante en la toma de decisiones estratégicas, priorizando:

\begin{itemize}
    \item La robustez del \textit{pipeline} modular sobre el volumen de procesamiento masivo.
    
    \item El desarrollo de un \textit{Test Set} representativo (de 2 a 3 páginas críticas por idioma) que permitiera validar el sistema de forma empírica sin necesidad de procesar los libros completos en esta fase inicial, optimizando el consumo de cuotas diarias de la API.
\end{itemize}

\subsection{Fases del desarrollo}

El ciclo de vida del software se ha dividido en cuatro fases principales, reflejadas en la bitácora de desarrollo y el repositorio:

\begin{itemize}
    \item \textbf{Fase de exploración y prototipado (Zero-Shot):} En las primeras etapas, se realizaron pruebas de extracción de texto plano para evaluar la capacidad de visión de Gemini. Se detectó la necesidad de transicionar hacia una salida estructurada debido a la pérdida de jerarquía semántica.
    
    \item \textbf{Fase de ingeniería de prompts y estructuración (Few-Shot):} Se implementó la lógica de \textit{Few-Shot Learning} (visible en la carpeta \texttt{data/few-shot}), proporcionando al modelo ejemplos de XML validado para guiar la generación. Metodológicamente, se optó por la técnica avanzada de pre-carga de historial (\textit{History Pre-loading}) a nivel de SDK, inyectando pares de documento de entrada y XML directamente en el búfer de la sesión antes del envío del mensaje del usuario, lo que optimizó drásticamente la latencia y la eficiencia en la transferencia de red. En esta fase se definió el DTD inicial.
    
    \item \textbf{Fase de refinamiento estructural y validación:} Desarrollo de los módulos de post-procesamiento (\texttt{xml\_corrector.py} y \texttt{xml\_validator.py}). Se iteró sobre el código para subsanar los errores estructurales típicos del modelo (etiquetas mal cerradas o duplicadas). En esta fase se diseñó la lógica de escape del bucle de validación (\textit{Validation Loop}): si tras el quinto intento de autoreparación iterativa el XML generado persiste en incumplir las especificaciones gramaticales del archivo DTD, el orquestador interrumpe las llamadas para evitar el bloqueo del \textit{pipeline}, guardando la última traza con errores heurísticos y continuando con el flujo secuencial de ejecución.
    
    \item \textbf{Fase de experimentación y análisis:} Realización de pruebas sistemáticas variando la temperatura del modelo (de 0.0 a 2.0) para identificar el ``punto dulce'' de precisión, tal como se detalla en el análisis de resultados del Capítulo~\ref{cap:resultados}. Durante esta etapa (específicamente en los ensayos del 10 de febrero de 2026), se descubrió una limitación física insalvable en el procesamiento por bloques: el sistema se encuentra acotado a un máximo de 2 páginas simultáneas. Debido a que cada sílaba analizada lingüísticamente en las glosas exige un mínimo de 4 etiquetas XML anidadas (aproximadamente 6 líneas de marcado), un volumen superior satura la ventana máxima de tokens de salida de la API, provocando truncamientos estructurales fatales.
\end{itemize}

\subsection{Control de versiones y documentación}

La trazabilidad del proyecto se ha garantizado mediante el uso de un repositorio en GitHub, organizando el código en una arquitectura modular que separa la lógica de procesamiento (\texttt{core/}), la configuración (\texttt{config/}), la validación (\texttt{validator/}) y la evaluación (\texttt{evaluator/}). Este ordenamiento facilita la mantenibilidad del sistema y su posible escalabilidad futura.
    

\section{Configuración del entorno de desarrollo}

La estabilidad de un sistema de Inteligencia Artificial depende críticamente de la gestión de dependencias y el aislamiento del entorno de ejecución. Para este proyecto, se ha diseñado un ecosistema basado en estándares de la industria que garantiza la portabilidad del código.

\subsection{Gestión de entornos virtuales (Anaconda)}

Se ha utilizado Anaconda como gestor de entornos y paquetes, permitiendo aislar las librerías del TFM del resto del sistema operativo. Esta decisión técnica evita conflictos entre versiones de Python y facilita la replicabilidad por parte de otros investigadores.

\begin{itemize}
    \item \textbf{Nombre del entorno:} \texttt{gemini-env}.
    
    \item \textbf{Versión de Python:} 3.12 (seleccionada por su optimización en la gestión de hilos y compatibilidad con las últimas librerías de Google AI).
    
    \item \textbf{Procedimiento de inicialización:} El entorno se activa mediante el comando \texttt{conda activate gemini-env}, asegurando que todas las variables de ruta apuntan a las binarias correctas.
\end{itemize}

\textbf{Nota de resolución de entorno:} Durante la fase inicial de despliegue, se documentó un conflicto de resolución de entorno derivado del uso del comando abreviado \texttt{py} en sistemas Windows. Dicho comando puenteaba el entorno virtual activo de Anaconda ejecutando el intérprete global del sistema, lo que impedía la localización de la SDK \texttt{google-genai}. El problema se solventó forzando metodológicamente la sintaxis explícita \texttt{python main.py}, garantizando así la vinculación estricta de las variables de entorno de \texttt{gemini-env}.

\subsection{Dependencias y librerías críticas}

El sistema se apoya en un conjunto de librerías especializadas instaladas a través de \texttt{pip} y registradas en el archivo \texttt{requirements.txt}. Las dependencias clave se agrupan en tres categorías:

\paragraph{Interacción con modelos de IA}

\begin{itemize}
    \item \textbf{\texttt{google-genai}:} Librería principal empleada para la comunicación con la API de Gemini. En el proyecto se utiliza para crear el cliente de inferencia, enviar documentos PDF en formato binario mediante \texttt{types.Part.from\_bytes} y configurar parámetros de generación como el modelo y la temperatura. También se emplea en el módulo de corrección automática de XML, donde se realiza una segunda llamada al modelo para corregir salidas que no cumplen el DTD.
\end{itemize}

\paragraph{Procesamiento de documentos PDF}

\begin{itemize}
    \item \textbf{\texttt{PyPDF2}:} Librería utilizada para leer documentos PDF, calcular el número total de páginas y extraer rangos concretos del documento. El módulo \texttt{pdf\_processor.py} emplea \texttt{PdfReader} y \texttt{PdfWriter} para generar un nuevo PDF en memoria, almacenado como un flujo de bytes mediante \texttt{io.BytesIO}. Este flujo binario es el que posteriormente se envía a Gemini como entrada documental.
\end{itemize}

\paragraph{Validación y corrección de XML}

\begin{itemize}
    \item \textbf{\texttt{lxml}:} Librería utilizada en el módulo \texttt{xml\_validator.py} para comprobar que la salida XML generada por la IA está bien formada y cumple las reglas definidas en el archivo \texttt{gramatica.dtd}. Para ello, el sistema parsea el XML, carga el DTD y valida la estructura antes de aceptar la salida como válida.

    \item \textbf{\texttt{re}:} Librería estándar de Python empleada para limpiar la respuesta generada por el modelo, eliminando posibles delimitadores de Markdown como \texttt{```xml} o \texttt{```} antes de procesar el contenido como XML.
\end{itemize}

\paragraph{Gestión de ejecución, configuración y registros}

\begin{itemize}
    \item \textbf{\texttt{os}:} Librería estándar utilizada para comprobar la existencia de archivos, gestionar rutas y crear directorios necesarios durante la ejecución, como el directorio de logs.

    \item \textbf{\texttt{logging}:} Librería estándar empleada para registrar el flujo de ejecución del sistema. El módulo \texttt{logger\_config.py} configura un logger global que escribe tanto en consola como en archivos de log con marca temporal.

    \item \textbf{\texttt{time}:} Librería estándar utilizada para introducir esperas entre reintentos cuando se producen errores temporales en la API, como sobrecarga del servidor o fallos 503.

    \item \textbf{\texttt{datetime}:} Librería estándar utilizada para generar nombres de archivo de log con fecha y hora, facilitando la trazabilidad de las ejecuciones.
\end{itemize}
    

\subsection{Estructura de directorios}

El entorno se ha organizado siguiendo una arquitectura limpia (\textit{Clean Architecture}), donde el script principal \texttt{main.py} actúa como orquestador, invocando módulos específicos para cada tarea. Esta modularidad permite, por ejemplo, actualizar el motor de conversión a \LaTeX{} (\texttt{converter.py}) sin afectar a la lógica de extracción de la IA (\texttt{ai\_generator.py}).


\section{Gestión de claves de API y seguridad de la información}

En proyectos que dependen de servicios en la nube facturables o sujetos a cuotas restrictivas, la seguridad de las credenciales de acceso es una prioridad crítica. Durante el desarrollo de este TFM, se han implementado protocolos para asegurar que la API Key de Gemini sea gestionada de forma externa al código fuente.

\subsection{Externalización de credenciales}

Siguiendo los estándares de seguridad de software, se ha evitado la práctica de ``hardcodear'' (insertar directamente) las claves en los scripts de Python. En su lugar, se ha optado por un sistema de variables de entorno:

\begin{itemize}
    \item \textbf{Inyección en tiempo de ejecución:} La clave se carga en la memoria del sistema operativo antes de ejecutar la aplicación mediante el comando: \\ \texttt{set GEMINI\_API\_KEY=[CLAVE\_SECRETA]}.
    
    \item \textbf{Consumo mediante código:} El módulo \texttt{config/properties.py} utiliza la librería \texttt{os} para recuperar esta clave de forma dinámica, asegurando que el valor real nunca quede registrado en el historial de archivos del proyecto.
\end{itemize}

\subsection{Protección del repositorio y prevención de fugas}

Para garantizar que la información sensible no sea publicada accidentalmente en el repositorio de GitHub, se han aplicado las siguientes medidas:

\begin{itemize}
    \item \textbf{Uso de \texttt{.gitignore}:} Se ha configurado un archivo de exclusión que impide que archivos de configuración local (como \texttt{.env}), registros de logs con trazas de depuración o carpetas de entorno virtual sean subidos al servidor público.
    
    \item \textbf{Aislamiento de logs:} El sistema de registro implementado en \texttt{core/logger\_config.py} ha sido configurado para monitorizar el flujo de la aplicación sin volcar en los archivos de texto (\texttt{logs/*.log}) el contenido íntegro de las peticiones de API que contienen la clave de acceso.
\end{itemize}

\subsection{Integridad de los datos de la ALMG}

En cumplimiento con las restricciones de propiedad intelectual mencionadas en el Capítulo 1, el entorno de desarrollo se ha configurado para que los archivos PDF originales de la Academia de Lenguas Mayas de Guatemala residan únicamente en el almacenamiento local del investigador. El repositorio en la nube solo contiene los scripts de procesamiento y los resultados validados, garantizando el respeto a los derechos de autor de las fuentes originales.

\section{Selección y descripción del corpus de trabajo}

La validez de este estudio depende de la representatividad y complejidad del material documental seleccionado. Para ello, se ha conformado un corpus de trabajo compuesto por cuatro gramáticas de lenguas mayas, editadas por la Academia de Lenguas Mayas de Guatemala (ALMG), que representan diferentes variantes y desafíos técnicos: Kaqchikel, K’iche’, Mam y Q’anjob’al.

\subsection{Criterios de selección}

Se han seleccionado estas obras basándose en tres criterios estratégicos:

\begin{itemize}
        \item \textbf{Relevancia demográfica y representatividad documental:} El K’iche’, el Kaqchikel, el Mam y el Q’anjob’al permiten evaluar el sistema sobre lenguas mayas con distinta disponibilidad documental, variedad tipográfica y complejidad estructural, lo que refuerza la representatividad del corpus utilizado.
    
    \item \textbf{Diversidad estructural:} Cada gramática presenta una disposición tipográfica distinta, lo que permite evaluar la resiliencia del sistema ante diferentes estilos editoriales.
    
    \item \textbf{Complejidad del soporte:} El corpus incluye tanto documentos digitales con capa de texto como archivos consistentes en fotografías de libros físicos, permitiendo testear las capacidades multimodales del modelo Gemini.
\end{itemize}

\subsection{Caracterización del material documental}

A continuación, se detallan las fuentes que integran el corpus de estudio:

\begin{itemize}
    \item \textbf{Gramática Normativa Kaqchikel:} Representa el mayor desafío estructural del proyecto. Se ha identificado como una obra con alta densidad de glosas interlineales y tablas de paradigmas verbales complejos. Especialmente, la página 172 de este documento se ha utilizado como ``patrón de estrés'' para las pruebas de temperatura debido a su saturación informativa.
    
    \item \textbf{Gramática Normativa K'iche':} Este documento destaca por su rigor en el análisis morfológico. Ha servido para validar la capacidad de la IA en la extracción de segmentaciones lingüísticas y su correcta traducción al castellano en bloques adyacentes.
    
    \item \textbf{Gramática Normativa Mam:} Incorporada para asegurar que el pipeline sea capaz de manejar variaciones en el conjunto de caracteres Unicode y estructuras de numeración jerárquica que difieren levemente de las otras variantes.

    \item \textbf{Gramática Descriptiva Q’anjob’al:} Incluida para ampliar la validación del sistema a una cuarta lengua maya y comprobar su comportamiento ante una obra con características documentales y tipográficas distintas. Sus páginas 44, 94 y 201 forman parte del conjunto experimental empleado en la evaluación.
\end{itemize}

\subsection{Conformación del test set (Conjunto de pruebas)}

Dada la restricción de tiempo y potencia de cómputo (analizada en el Capítulo 1), no se ha procedido al procesado lineal de los libros completos. En su lugar, se ha construido un \textit{Test Set} compuesto por una selección de 2 a 3 páginas de alta complejidad por cada idioma.

Este subconjunto de datos ha sido etiquetado manualmente para generar el \textit{Ground Truth} necesario, permitiendo que el módulo \texttt{evaluator} realice comparativas estadísticas precisas sobre la estabilidad del sistema bajo diferentes configuraciones de temperatura.











% \chapter{Metodología} \label{metodologia}



% \section{Marco de Trabajo: Scrum Adaptado}
% Durante el desarrollo de SIDManager, adoptamos una metodología ágil basada en Scrum, adaptada a las necesidades del proyecto y a la colaboración remota con Santander Consumer Finance. Implementamos:

% \begin{itemize}
%     \item Sprints Semanales: Cada lunes, realizábamos una planificación de sprint con el tutor de Santander, donde definíamos los objetivos clave y priorizábamos tareas usando diagramas UML de secuencia.
%     \item Daily Standups Virtuales: Breves reuniones matutinas (15 min) para sincronizar avances y bloquedores, manteniendo al equipo alineado a pesar de la distancia.
%     \item Retrospectivas: Al final de cada sprint, analizábamos qué funcionó (ej: code reviews estructuradas) y qué mejorar (ej: ajustar la granularidad de las user stories).
% \end{itemize}

% \section{Herramientas Ágiles Clave}

% Para garantizar trazabilidad y calidad:

% \begin{itemize}
%     \item Code Reviews Colaborativas: Se realizaron revisiones de código semanales, donde el tutor evaluaba los commits contra estándares como Java Code Conventions y patrones de diseño como Singleton o Facade.
%     \item UML como Lenguaje Común: Diagramas de secuencia  servían como esquemas técnicos antes de implementar features críticas (ej: flujo de autenticación con Spring Security).
%     \item Tablero Kanban Digital: En Trello, gestionábamos tareas con columnas clásicas (To Do, In Progress, Done) y etiquetas de prioridad.
% \end{itemize}

% \section{Adaptación a Cambios}
% La naturaleza iterativa de Scrum nos permitió:

% \begin{itemize}
%     \item Incorporar Feedback en Tiempo Real: Cuando se identificó la necesidad de pseudopaginación para 20K documentos, ajustamos el sprint siguiente para priorizarlo.
%     \item Mantener Calidad bajo Presión: Las definition of done incluían: (1) pruebas manuales, (2) logs estructurados con Log4j, y (3) documentación Javadoc.
% \end{itemize}

% Este enfoque en metodologías ágiles generó resultados tangibles.

% \begin{itemize}
%     \item Entregables Consistentes: 8 sprints completados, con demo funcional tras cada uno.
%     \item Código Limpio y Documentado: +15 clases con Javadoc, facilitando la transición al equipo de Santander.
%     \item Flexibilidad: Algunos cambios mayores de requisitos absorbidos sin impactar la fecha de entrega (como la migración del inicio de sesión manual a Spring Security).
% \end{itemize}


